\chapter{NKMemory : qui alloue, qui libère}

C'est le chapitre le plus important des fondations, et le seul dont une erreur ne
se manifeste jamais là où elle a été commise.

Nkentseu n'emploie ni \texttt{new}, ni \texttt{delete}, ni \texttt{malloc}. Toute
la mémoire passe par \textbf{NKMemory}. Ce n'est pas une préférence de style :
c'est ce qui rend possible le suivi des fuites, les stratégies d'allocation par
usage, et le portage sur des plateformes où le tas du système n'est pas celui
qu'on croit.

\section{La règle, et ce qu'elle coûte de violer}

\begin{nkretenir}
\textbf{Ce qui a été alloué par NKMemory se libère par NKMemory. Point.}

Un bloc obtenu par \texttt{NkAlloc} et rendu par \texttt{free} --- ou par
\texttt{delete[]} --- rend un pointeur à un gestionnaire de tas \emph{qui ne l'a
jamais donné}. Sous Windows, cela produit le code \texttt{0xC0000374},
« corruption du tas ».

Et voici pourquoi c'est le pire des défauts : \textbf{le programme ne plante pas
au moment de la faute}. Il plante plus tard, dans une allocation parfaitement
innocente, parfois des milliers d'instructions après --- et c'est cette
allocation-là qu'on accuse.
\end{nkretenir}

\begin{nkpiege}[Trois cas réels du dépôt]
\begin{enumerate}
\item Le tampon rendu par un encodeur --- \texttt{NkJPEGCodec::Encode},
      \texttt{NkPNGCodec::Encode}, \texttt{SaveToMemory} --- est alloué par
      \texttt{NkAlloc}. Le libérer avec \texttt{delete[]} corrompt le tas.
\item \texttt{NkImage::Alloc()} suivi de \texttt{Free()} \emph{puis} de
      \texttt{delete img} est une double libération : \texttt{Free()} rend déjà
      l'enveloppe.
\item Appeler \texttt{Free()} sur une \texttt{NkImage} déclarée \textbf{sur la
      pile} exécute une libération \emph{sur une adresse de pile}. Cela a
      réellement fait planter une démo du dépôt.
\end{enumerate}
\textbf{La question à se poser devant toute libération n'est pas « comment
libère-t-on ceci ? » mais « qui l'a alloué ? ».}
\end{nkpiege}

\section{Les quatre fonctions de base}

\begin{nkcode}{Kernel/Foundation/NKMemory/src/NKMemory/NkAllocator.h:704-737}
void *NkAlloc(nk_size size, NkAllocator *allocator = nullptr,
              nk_size alignment = NK_MEMORY_DEFAULT_ALIGNMENT);

void *NkAllocZero(nk_size count, nk_size size, NkAllocator *allocator = nullptr,
                  nk_size alignment = NK_MEMORY_DEFAULT_ALIGNMENT);

void *NkRealloc(void *ptr, nk_size oldSize, nk_size newSize,
                NkAllocator *allocator = nullptr,
                nk_size alignment = NK_MEMORY_DEFAULT_ALIGNMENT);

void NkFree(void *ptr, NkAllocator *allocator = nullptr);
void NkFree(void *ptr, nk_size size, NkAllocator *allocator = nullptr);
\end{nkcode}

\begin{nkretenir}
\textbf{Le paramètre \texttt{allocator} à \texttt{nullptr} signifie « celui par
défaut »}, pas « aucun ». C'est ce qui rend le code ordinaire lisible ---
\texttt{NkAlloc(1024)} suffit --- tout en permettant à un sous-système d'imposer
le sien.

\textbf{Et il doit être le MÊME à l'allocation et à la libération.} Allouer sur
un allocateur nommé et libérer sur celui par défaut est exactement la même faute
que mélanger \texttt{NkAlloc} et \texttt{free}, avec un déguisement de plus.
\end{nkretenir}

\begin{nkretenir}
\textbf{\texttt{NkRealloc} rend \texttt{nullptr} en cas d'échec --- et l'ancien
pointeur reste valide.} C'est écrit dans sa documentation, et c'est la seule
façon correcte : écraser votre pointeur par le retour sans le tester perd le
bloc d'origine.

\begin{nkcode}{FAUX}
p = NkRealloc(p, ancien, neuf);   // si echec : p vaut nullptr, et le bloc est perdu
\end{nkcode}
\end{nkretenir}

\section{Allouer un objet, pas des octets}

\texttt{NkAlloc} rend de la mémoire brute : aucun constructeur n'a été appelé.
Pour un objet, ce sont les quatre modèles de \texttt{NkAllocator} qu'il faut :

\begin{nkcode}{Kernel/Foundation/NKMemory/src/NKMemory/NkAllocator.h:239-273}
template <typename T, typename... Args> [[nodiscard]] T *New(Args &&...args);
template <typename T> void Delete(T *ptr) noexcept;
template <typename T, typename... Args> [[nodiscard]] T *NewArray(SizeType count, Args &&...args);
template <typename T> void DeleteArray(T *ptr) noexcept;
\end{nkcode}

\begin{nkretenir}
\textbf{Chacun fait deux choses, et les fait dans le bon ordre.}
\texttt{New<T>} alloue \emph{puis} construit ; \texttt{Delete<T>} détruit
\emph{puis} libère. C'est exactement ce que font \texttt{new} et \texttt{delete}
du langage --- mais sur \emph{votre} allocateur.

Et \texttt{NewArray}/\texttt{DeleteArray} vont ensemble : le premier range le
nombre d'éléments à côté du tableau, le second le relit pour appeler le bon
nombre de destructeurs. Libérer un tableau avec \texttt{Delete<T>} n'appelle
qu'un seul destructeur et se trompe d'adresse de base.
\end{nkretenir}

\begin{nkpiege}
\texttt{[[nodiscard]]} sur \texttt{New} n'est pas décoratif : ignorer le retour
d'une allocation est une fuite garantie, et le compilateur le refuse désormais.
C'est le genre de garde qui coûte un mot et travaille pour toujours.
\end{nkpiege}

\section{Neuf allocateurs, et pourquoi il en faut plusieurs}

Un tas généraliste doit répondre à n'importe quelle demande dans n'importe quel
ordre. C'est un problème difficile --- donc lent. Or la plupart des programmes
n'ont pas ce besoin : ils allouent \emph{par motif}.

\begin{center}
\begin{tabular}{@{}lll@{}}
\toprule
\textbf{Allocateur} & \textbf{Motif qu'il exploite} & \textbf{Ce qu'il abandonne} \\
\midrule
\texttt{NkMallocAllocator} & aucun --- le tas du système & rien, mais rien de gagné \\
\texttt{NkNewAllocator} & aucun --- \texttt{new}/\texttt{delete} & idem \\
\texttt{NkVirtualAllocator} & très gros blocs, pages OS & la granularité fine \\
\texttt{NkLinearAllocator} & « tout meurt en même temps » & la libération individuelle \\
\texttt{NkArenaAllocator} & idem, par zones & idem \\
\texttt{NkStackAllocator} & dernier alloué, premier libéré & l'ordre libre \\
\texttt{NkPoolAllocator} & objets de \emph{même taille} & les tailles variables \\
\texttt{NkFreeListAllocator} & tailles variables, réemploi & la simplicité \\
\texttt{NkBuddyAllocator} & tailles en puissances de deux & jusqu'à 50\,\% de place \\
\bottomrule
\end{tabular}
\end{center}

\begin{nkretenir}
\textbf{Chaque ligne de la colonne du milieu est une \emph{hypothèse sur la durée
de vie}}, et c'est elle qu'on achète. Un allocateur rapide n'est pas rapide en
soi : il est rapide \emph{parce qu'il refuse de répondre à des questions que
vous ne poserez pas}.

Le corollaire est immédiat : \textbf{choisir un allocateur, c'est déclarer un
motif d'usage}. Si le motif est faux, l'allocateur est plus lent que le tas
ordinaire, ou il échoue.
\end{nkretenir}

\subsection{Le cas le plus utile : l'allocateur linéaire}

\begin{nkcode}{Kernel/Foundation/NKMemory/src/NKMemory/NkAllocator.h:393-398}
// Tres rapide pour les allocations temporaires frame-based.
// Ne supporte pas la deallocation individuelle - Reset() libere tout.
\end{nkcode}

\begin{nkretenir}
Un allocateur linéaire n'est qu'un \textbf{curseur qui avance}. Allouer, c'est
ajouter la taille au curseur et rendre l'adresse précédente : quelques
instructions, aucune recherche, aucune fragmentation possible.

\texttt{Deallocate} n'y fait presque rien --- l'en-tête le dit : il n'optimise
que le cas où l'on rend le \emph{dernier} bloc. La libération réelle est
\texttt{Reset()}, qui remet le curseur à zéro et rend tout d'un coup.

\textbf{Le patron d'emploi} : un allocateur linéaire par image. Tout ce qui est
temporaire y est alloué ; à la fin de l'image, un \texttt{Reset()} et le coût de
libération de milliers d'objets tombe à une affectation.
\end{nkretenir}

\begin{nkpiege}
\textbf{Rien de ce qui est alloué sur un allocateur linéaire ne doit survivre à
son \texttt{Reset()}.} Garder un pointeur d'une image sur l'autre donne une
adresse qui pointe désormais sur \emph{autre chose} --- pas sur du vide, ce qui
serait détectable, mais sur des données valides appartenant à quelqu'un d'autre.

C'est le pointeur pendant sous sa forme la plus perverse : la mémoire est
lisible, le programme ne plante pas, et les valeurs sont fausses.
\end{nkpiege}

\begin{nkretenir}
\texttt{Capacity()}, \texttt{Used()} et \texttt{Available()} sont là pour être
\emph{lus}. Un allocateur linéaire qui déborde ne s'agrandit pas : il échoue. La
bonne pratique est de journaliser le pic d'utilisation en développement, et de
dimensionner sur la mesure --- jamais sur une estimation.
\end{nkretenir}

\section{Savoir ce qu'on a alloué}

\begin{center}
\begin{tabular}{@{}ll@{}}
\toprule
\textbf{Outil} & \textbf{Ce qu'il apporte} \\
\midrule
\texttt{NkTracker} & détection de fuites en $O(1)$, par table de hachage \\
\texttt{NkTag} & étiqueter une allocation --- « textures », « audio »\dots \\
\texttt{NkProfiler} & compteurs, pics, statistiques par allocateur \\
\texttt{NkGc} & un ramasse-miettes \emph{marquage et balayage}, pour qui le veut \\
\bottomrule
\end{tabular}
\end{center}

\begin{nkretenir}
\textbf{Le suivi n'est possible que parce que tout passe par NKMemory.} C'est le
retour sur investissement de la règle du début de chapitre : un seul \texttt{new}
échappé, et le bloc correspondant devient invisible au traqueur.

Autrement dit, la discipline n'est pas une contrainte \emph{en échange de} rien :
elle est ce qui achète l'instrument.
\end{nkretenir}

\begin{nkpiege}
\textbf{Un traqueur de fuites signale ce qui n'a pas été libéré \emph{à la fin du
programme}.} Il ne voit pas une mémoire qui grossit indéfiniment mais qui finit
par être rendue --- et c'est pourtant le défaut qui tue une application longue.

Pour celui-là, l'instrument est le \emph{pic} d'utilisation dans le temps, pas le
solde à la fermeture. Deux mesures différentes, deux défauts différents.
\end{nkpiege}

\section{Ce que vous écrirez vraiment}

En pratique, dans du code d'application, vous n'appellerez presque jamais
\texttt{NkAlloc} :

\begin{center}
\begin{tabular}{@{}ll@{}}
\toprule
\textbf{Besoin} & \textbf{Ce qu'on emploie} \\
\midrule
une suite d'éléments & \texttt{NkVector} --- il alloue pour vous \\
du texte & \texttt{NkString} \\
un objet qui vit tant qu'on le tient & \texttt{NkUniquePtr} \\
un tampon brut, une image, un fichier & \texttt{NkAlloc}/\texttt{NkFree} \\
\bottomrule
\end{tabular}
\end{center}

\begin{nkretenir}
\textbf{Les conteneurs et les pointeurs intelligents \emph{sont} la façon normale
d'employer NKMemory.} Ils allouent à travers lui, et ils libèrent tout seuls.

Ce chapitre n'est donc pas un mode d'emploi quotidien : c'est ce qu'il faut avoir
lu pour comprendre ce que le conteneur fait à votre place --- et pour savoir quoi
regarder le jour où le tas se corrompt.
\end{nkretenir}

\begin{nkbilan}
Toute la mémoire de Nkentseu passe par NKMemory, et \textbf{ce qui a été alloué
par lui se libère par lui} : mélanger les gestionnaires corrompt le tas, et le
plantage arrive plus tard, ailleurs, dans du code innocent. \texttt{NkAlloc} rend
des octets, \texttt{New<T>} rend un objet construit --- et \texttt{NewArray} va
avec \texttt{DeleteArray}, jamais avec \texttt{Delete}. Les neuf allocateurs
n'existent pas pour le choix : chacun \emph{achète de la vitesse en refusant de
répondre à une question}, et le linéaire --- allouer sans jamais libérer
individuellement, \texttt{Reset()} à la fin de l'image --- est celui qui rapporte
le plus. Enfin, le suivi des fuites n'est possible que parce que la règle est
tenue partout.
\end{nkbilan}

\begin{nkquestions}
\begin{enumerate}
\item Que se passe-t-il si l'on libère avec \texttt{free} un bloc obtenu par
      \texttt{NkAlloc} --- et \emph{quand} cela se voit-il ?
\item Quelle différence entre \texttt{NkAlloc} et \texttt{New<T>} ?
\item Pourquoi \texttt{NewArray} exige-t-il \texttt{DeleteArray} ?
\item Que fait \texttt{Deallocate} sur un allocateur linéaire, et comment y
      libère-t-on réellement ?
\item Pourquoi un traqueur de fuites ne voit-il pas une mémoire qui grossit puis
      est rendue ?
\end{enumerate}
\end{nkquestions}

\begin{nkpratique}
Écrivez un \texttt{NkLinearAllocator} de 64 Kio et servez-vous-en pour allouer,
à chaque image d'une petite application, un tableau temporaire dont la taille
dépend de ce qui est à l'écran.

Trois exigences :
\begin{enumerate}
\item afficher \texttt{Used()} et \texttt{Available()} en clair, à chaque image ;
\item appeler \texttt{Reset()} à la fin de l'image, et \emph{nulle part
      ailleurs} ;
\item faire délibérément déborder l'allocateur, et voir ce qui se passe. Notez
      ce que l'allocateur rend, et ce que votre code en fait --- c'est la
      deuxième partie qui compte.
\end{enumerate}
\end{nkpratique}

\begin{nkdemo}
Prouvez qu'un pointeur survivant à un \texttt{Reset()} donne des valeurs fausses
sans planter.

Allouez une structure sur l'allocateur linéaire à l'image $n$, remplissez-la avec
des valeurs que vous reconnaîtrez, et \textbf{gardez le pointeur}. À l'image
$n+1$, après le \texttt{Reset()} et après avoir alloué autre chose, relisez la
structure et affichez son contenu.

\textbf{Ce que la démonstration doit établir}, et c'est le point : non pas que le
programme plante --- il ne plantera pas --- mais que les valeurs relues sont
\emph{plausibles et fausses}. Affichez-les à côté des valeurs d'origine.

\alerte{} Ne concluez rien d'une seule exécution où les valeurs n'auraient pas changé :
cela signifie seulement que rien n'a encore réécrit cette zone. Allouez
davantage à l'image $n+1$ et recommencez.
\end{nkdemo}

\begin{nklogique}
Une application plante avec \texttt{0xC0000374} au bout de dix minutes, toujours
dans une allocation différente, jamais au même endroit.

\begin{enumerate}
\item Expliquez pourquoi le point de plantage ne désigne pas le coupable.
\item Énumérez au moins quatre gestes du code susceptibles de produire ce
      symptôme, et rangez-les du plus au moins probable.
\item Vous soupçonnez un module précis. Décrivez comment vous le
      \emph{disculperiez} --- pas comment vous le confirmeriez. Expliquez
      pourquoi c'est le bon sens de l'enquête ici, alors que la démarche
      inverse serait naturelle.
\end{enumerate}
\end{nklogique}
